MQOM (``MQ on my Mind'') is a post-quantum digital signature scheme, built on the
Multi-Party-Computation-in-the-Head (MPCitH) / Threshold-Computation-in-the-Head (TCitH)
paradigm, and advanced to the third round of the NIST Additional Digital Signatures
process. While MQOM's \emph{mathematical} security rests on the hardness of the
Multivariate Quadratic (MQ) problem, its \emph{implementation} security against
side-channel analysis is comparatively unexplored. This project investigates whether
the permanent secret witness $x$ of MQOM~v2.1 can be recovered by Correlation Power
Analysis (CPA) acting on power leakage of secret-dependent intermediate values.

The study is conducted entirely through mathematics and computer-based simulation. We
do not use physical hardware, measurement equipment, real power or electromagnetic
traces, or the official MQOM reference implementation. Instead, we build independent
Python implementations of the relevant finite-field operations and MQOM intermediate
values derived from the published specification, generate synthetic leakage under
explicitly stated hypothetical leakage models, and apply CPA-based distinguishers. We
are careful throughout to separate specification-established facts, published results,
our own mathematical deductions, the assumptions used to construct a simulation, and
the results obtained from that simulation. In particular, synthetic leakage is never
presented as a physical measurement, and a successful simulated recovery is never
presented as a demonstrated vulnerability of a real MQOM device.

From an initial attack surface of seven candidate leakage locations we isolate those
that expose the secret key to \emph{known, per-signature-varying} public data, a
precondition for CPA. We decompose $x$ into a prefix, consumed by the correlated seed
tree, and a suffix, whose correction coordinates are public in the signature. For the
suffix we formulate a vertical key-XOR CPA. Our simulations show that a na\"ive
absolute-correlation distinguisher suffers an irreducible two-way (complement)
ambiguity intrinsic to any affine leakage on a bare-XOR target, and that a signed
distinguisher aligned with the assumed leakage polarity resolves it, recovering the
full suffix under the stated model. The project further studies robustness to
leakage-model mismatch and the recovery of the prefix, working towards reconstruction
of the complete witness, together with an honest account of what simulation can and
cannot establish.
